\section{Preliminaries}
\label{sec:preliminary}

This section reviews the attention computation underlying autoregressive language models and block sparse attention, which reduces long-context inference cost by restricting each query to a subset of context blocks.

\paragraph{Standard attention.}
Let $\mathbf{q} \in \mathbb{R}^{d}$ denote a query vector of the current token, and $\mathbf{K}, \mathbf{V} \in \mathbb{R}^{n \times d}$ denote the key and value matrices of the $n$ previous tokens, which are referred to as \textit{context tokens}.
Standard attention \cite{vaswani2017attention} computes the output $\mathbf{o} \in \mathbb{R}^d$ as
\begin{equation}
    \mathbf{o} = \operatorname{softmax}(\mathbf{q} \mathbf{K}^\top) \mathbf{V}.
    \label{eqn:fula}
\end{equation}
Here we omit the scaling factor for notational simplicity. During autoregressive decoding, each new query attends to all previous keys and values. Since the number of previous tokens grows linearly with the decoding step, the cumulative attention cost over a sequence of length $n$ grows as $\mathcal{O}(n^2)$.

\paragraph{Block sparse attention.}
To mitigate this quadratic cost, block sparse attention \cite{zhangsurvey,sun2025efficient} restricts the query $\mathbf{q}$ to attend only to a selected subset of context blocks.
Let $\mathcal{S}\subseteq\{1,\ldots,n\}$ denote the selected token-index set. 
The attention output is formulated as
\begin{equation}
    \mathbf{o}
    =
    \operatorname{softmax}
    \left(
        \mathbf{q}\mathbf{K}_{\mathcal{S}}^\top
    \right)
    \mathbf{V}_{\mathcal{S}}.
    \label{eqn:spa}
\end{equation}
For a single query, this reduces the computation cost from $\mathcal{O}(n)$ to $\mathcal{O}(|\mathcal{S}|)$.
Over a sequence of length $n$, the cumulative cost is therefore reduced to $\mathcal{O}(n|\mathcal{S}|)$.
To construct $\mathcal{S}$, the $n$ context positions are partitioned into $C$ contiguous index blocks $\{B_1,\ldots,B_C\}$, where each $B_m$ denotes a set of token indices and contains $b$ tokens (i.e., $C=n/b$).
A lightweight selector $\mathcal{R}_\theta(\cdot)$ computes relevance scores $\mathbf{s}\in\mathbb{R}^{C}$ over these blocks.
The selector chooses the $K$ most important blocks from all $C$ context blocks, and the query only attends to the union of these selected blocks.
Formally, the selected block indices and the corresponding token-index set are
\begin{equation}
    \mathcal{I}
    =
    \text{Top-}K(\mathbf{s}, K),
    \qquad
    \mathcal{S}
    =
    \bigcup_{m\in\mathcal{I}} B_m .
    \label{eqn:selected_context}
\end{equation}
Under plain hard Top-$K$, the selected index set is piecewise constant with respect to selector scores. Therefore, the language modeling loss provides no useful gradient through the selection path to update the selector, as illustrated in \Cref{fig:motivation}a.
To make the selector end-to-end trainable, the core problem is how to make selector scores affect attention differentiably, so that the language modeling loss can train the selector.
